\chapter*{ЗАКЛЮЧЕНИЕ}
\addcontentsline{toc}{chapter}{ЗАКЛЮЧЕНИЕ}



В аналитической части работы был проведен анализ существующих методов конвертации двумерных изображений в трехмерные представления.
Методы были классифицированы по ключевым критериям: способу получения данных, типу исходных изображений, необходимости в дополнительной информации, результатам реконструкции, масштабу сцены и уровню автоматизации.

В ходе исследования были изучены различные подходы к восстановлению 3D-информации,
включая методы структурированного освещения, стерео сопоставления, DFD, Shape from Shading, Comfortable Depth Generation и нейросетевые методы, такие как MiDaS.

Проведенный анализ привел к выводу, что оптимальным решением для восстановления глубины из архивных фотографий является гибридный метод,
объединяющий преимущества нейросетевого подхода (MiDaS) и классического метода Shape from Shading.
Такое решение позволит обеспечить как общую точность восстановления геометрии сцены, так и детализацию отдельных элементов, сохраняя при этом устойчивость к дефектам, характерным для исторических изображений.

Представленная формальная модель задачи и декомпозиция метода на функциональные блоки предобработки, гибридного восстановления и постобработки формируют теоретическую основу для разработки практического решения в следующих разделах работы.



В конструкторском разделе была представлена детальная структура разработанного гибридного метода восстановления трехмерной информации из архивных фотографий.
Разработка включает семь основных алгоритмических блоков, последовательное выполнение которых обеспечивает преобразование входного изображения в детализированную карту глубины.

Представленные схемы алгоритмов являются основой для последующей программной реализации метода, которая будет рассмотрена в технологическом разделе работы.

В технологическом разделе выполнена полная программная реализация гибридного метода восстановления глубины из монокулярных изображений.

Для решения поставленной задачи был выбран язык Python с экосистемой специализированных библиотек (OpenCV, PyTorch, NumPy и др.),
что обеспечило эффективную работу с изображениями, применение моделей машинного обучения и визуализацию трехмерных данных.

Разработанный алгоритм успешно прошел тестирование на различных типах изображений (цветных, черно-белых, с эффектом сепии)
и продемонстрировал корректную обработку исключительных ситуаций, что подтверждает работоспособность созданного решения.



В исследовательском разделе проведено исследование с использованием стандартных датасетов,
алгоритмов искусственного старения изображений и объективных метрик оценки
для количественной оценки эффективности разработанного гибридного метода восстановления глубины.

Проведено исследование устойчивости гибридного метода к различным эффектам старения изображений,
которое подтвердило гипотезу о дифференцированном влиянии различных типов деградации на точность реконструкции глубины:
    \begin{itemize}
        \item эффекты, изменяющие структурную информацию (шум и виньетирование), оказывают наиболее сильное негативное влияние (26.8-42.4\% увеличения RMSE);
        \item эффекты, затрагивающие только цветовую гамму (сепия и выцветание), влияют минимально (0.26-1.6\% увеличения RMSE);
        \item локальные искажения (пятна и царапины) оказывают умеренное влияние (около 10.2\% увеличения RMSE).
    \end{itemize}

Выполнен сравнительный анализ разработанного гибридного метода с базовыми подходами (MiDaS и SFS) при различных уровнях деградации изображений:
    \begin{itemize}
        \item на качественных изображениях все три метода демонстрируют сопоставимые результаты по метрике RMSE;
        \item при увеличении уровня деградации метод SFS критически теряет в точности (RMSE возрастает на 45.4\%);
        \item MiDaS сохраняет стабильность на всех уровнях деградации;
        \item гибридный метод демонстрирует устойчивость на уровне MiDaS с улучшенными на 1-2\% абсолютными показателями точности.
    \end{itemize}

Результаты исследования подтверждают эффективность разработанного гибридного метода,
который успешно сочетает в себе устойчивость нейросетевого подхода к деградации изображений с детализирующими возможностями метода Shape from Shading.
Полученные данные позволяют определить оптимальные области применения метода:
восстановление трёхмерной информации из архивных фотографий для задач цифровой реставрации и трёхмерной реконструкции исторических объектов,
когда доступны только архивные фотоматериалы различного качества.
Метод позволяет получить более точные и детализированные карты глубины по сравнению с существующими подходами при работе с деградированными изображениями.


экономическую эффективность или значимость работы;
 прогнозы и предположения о возможных направлениях развития объекта
исследования.

